\documentclass[10pt]{beamer}

% Theme Choice
\usetheme{metropolis} 
\usepackage{appendixnumberbeamer}
\usepackage{booktabs}
\usepackage{amsmath, amsfonts, amssymb}

% Title Information
\title{Recommender? I Hardly Know Her!}
\subtitle{Policy-Guided Path Reasoning for Movie Recommendations}
\date{\today}
\author{Perry Santry, Joseph Crazylastname, Jesse Seeligsohn}
\institute{Comp 579 Project --- McGill University}

\begin{document}

\maketitle

\begin{frame}{Project Overview}
    \begin{itemize}
        \item Modern recommender systems often rely on "black-box" latent embeddings that lack interpretability.
        \item This project adapts \textbf{Policy-Guided Path Reasoning (PGPR)} to the \textbf{MovieLens-1M} dataset.
        \item Recommendation is treated as a \textbf{graph navigation problem} where an agent learns to find paths through a heterogeneous knowledge graph.
        \item The resulting paths serve as explicit explanations for the recommendations.
    \end{itemize}
\end{frame}

\section{Background}

\begin{frame}{Knowledge Graphs (KG)}
    \begin{itemize}
        \item Entities (users, movies, attributes) are connected by typed relations.
        \item The graph combines user-movie interactions with metadata (genre, director, etc.) from a Freebase-style schema.
        \item \textbf{Interpretable Reasoning Paths:}
        \begin{itemize}
            \item \textit{Content-based:} $user \to movie \to genre \to movie$.
            \item \textit{Collaborative:} $user \to movie \to user \to movie$.
        \end{itemize}
    \end{itemize}
\end{frame}

\section{Methodology}

\begin{frame}{State Encoding: A Design Choice}
    \begin{itemize}
        \item \textbf{Ekar Approach:} Uses an \textbf{LSTM} to encode the full sequence of entities and relations [cite: 208-210].
        \item \textbf{PGPR Approach (Our Implementation):} Uses \textbf{Concatenation} of the starting user, current entity, previous entity, and previous relation.
        \item \textbf{State Vector:} $s_t = [u; e_t; e_{t-1}; r_t]$.
        \item \textbf{Rationale:} Concatenation is computationally efficient and sufficient for short-horizon (3-hop) reasoning.
    \end{itemize}
\end{frame}

\begin{frame}{MDP and Soft Reward Strategy}
    \begin{itemize}
        \item \textbf{Horizon:} $T=3$ hops.
        \item \textbf{Soft Terminal Reward:} Scores relevance of the final movie $e_T$ to user $u$.
        \[ R_T = \begin{cases} \max(0, \frac{\langle u, e_T \rangle}{\max_{i \in I} \langle u, i \rangle}) & \text{if } e_T \in I \\ 0 & \text{otherwise} \end{cases} \] 
        \item Rewards are normalized to $[0, 1]$ to stabilize training across diverse user embedding scales.
    \end{itemize}
\end{frame}

\begin{frame}{Action Pruning and Policy Search}
    \begin{itemize}
        \item \textbf{User-Conditional Pruning:} Limits the action space to $\alpha = 400$ by scoring edges: $f((r, e)|u) = \langle u+r, e \rangle$.
        \item \textbf{Policy Network:} Two-layer MLP with ELU activations trained via REINFORCE.
        \item \textbf{Inference:} Beam search with $K=(25, 5, 1)$ to find diverse, high-probability recommendation paths.
    \end{itemize}
\end{frame}

\section{Analysis and Results}

\begin{frame}{Experimental Benchmarking (MovieLens-1M)}
    \begin{itemize}
        \item Comparison of reporting for KG-based methods on the MovieLens-1M dataset:
    \end{itemize}
    \begin{table}
        \centering
        \begin{tabular}{lrr}
            \toprule
            \textbf{Model} & \textbf{HR@10} & \textbf{NDCG@10} \\
            \midrule
            BPR-MF (Baseline) & 0.0895 & 0.1914 \\
            KTUP & 0.1579 & 0.3230 \\
            \textbf{Ekar (RL)} & \textbf{0.1889} & \textbf{0.3543} \\
            \bottomrule
        \end{tabular}
        \caption{Baseline results on MovieLens-1M.}
    \end{table}
\end{frame}

\begin{frame}{Frequent Path Patterns (Explainability)}
    \begin{itemize}
        \item Analysis of reasoning paths on MovieLens-1M reveals the agent's strategy:
        \item \textbf{Collaborative Logic (94.4\%):} $User \to Movie \to User \to Movie$.
        \item \textbf{Content Logic (1.5\%):} $User \to Movie \to Genre \to Movie$.
        \item \textbf{Finding:} The agent primarily discovers collaborative filtering paths, navigating to items liked by users with similar tastes.
    \end{itemize}
\end{frame}

\begin{frame}{Ablation Study: Why It Works}
    \begin{itemize}
        \item \textbf{Reward Shaping:} Crucial for convergence; without it, the agent receives zero reward for any item not already in the user's history.
        \item \textbf{Action Pruning:} Essential for filtering noise; smaller pruned spaces often lead to better performance by prioritizing high-relevance edges.
        \item \textbf{Path Length ($T=3$):} Increasing length (e.g., $T=5$) often degrades performance due to increased noise in long-order similarities.
    \end{itemize}
\end{frame}

\section{Conclusion}

\begin{frame}{Conclusion}
    \begin{itemize}
        \item Adapted the PGPR architecture to a movie recommendation setting.
        \item Demonstrated that RL agents can effectively navigate knowledge graphs to provide both accurate and interpretable results.
        \item Future work could involve incorporating more complex domain-specific reward functions.
    \end{itemize}
\end{frame}

\end{document}